\chapter{Joint Kinematics and Soft Tissue}
\label{ch:joint-kinematics}

\epigraph{%
  The knee is not a hinge. The shoulder is not a ball-and-socket.
  These are useful fictions that biology tolerates but does not enforce.
}{}

\section{Beyond Ideal Joints}

Engineering joints are designed to constrain motion to specific degrees of freedom:
a revolute joint permits one rotation, a prismatic joint permits one translation.
Biological joints achieve constraint through a fundamentally different mechanism:
the interaction of articular surfaces, ligaments, joint capsule, and muscle forces.
This section develops the kinematics of biological joints using the screw theory
framework of Volume~0, extended to handle the additional complexity.

\subsection{The Instantaneous Screw Axis}

Every rigid-body motion can be decomposed into a rotation about and translation along
a \textbf{screw axis} (Chasles' theorem, Volume~0, Chapter~5). For an engineering
revolute joint, this screw axis is fixed in both bodies. For a biological joint,
the instantaneous screw axis (ISA) \emph{moves} as the joint configuration changes.

\begin{definition}{Instantaneous Screw Axis of a Biological Joint}{isa-bio}
  Given the relative twist $\twist_{ij} = (\bm{\omega}, \bm{v}) \in \Reals^6$ between
  adjacent segments $i$ and $j$, the instantaneous screw axis has:
  \begin{itemize}
    \item Direction: $\hat{\bm{s}} = \bm{\omega} / \norm{\bm{\omega}}$
    \item Pitch: $h = \bm{\omega} \cdot \bm{v} / \norm{\bm{\omega}}^2$
    \item Location: $\bm{q} = \bm{\omega} \times \bm{v} / \norm{\bm{\omega}}^2$
  \end{itemize}
  For a pure rotation ($h = 0$), the ISA passes through $\bm{q}$ in direction $\hat{\bm{s}}$.
  Tracking how $\bm{q}$ and $\hat{\bm{s}}$ change with joint angle reveals the coupled
  kinematics of the biological joint.
\end{definition}

\subsection{The Knee: Rolling, Sliding, and Coupled Rotation}

The tibiofemoral joint exemplifies the complexity of biological kinematics. During
flexion from $0^\circ$ to $90^\circ$:
\begin{enumerate}
  \item The femoral condyles \textbf{roll} on the tibial plateaus (like a wheel on a road)
  \item They simultaneously \textbf{slide} posteriorly (the contact point moves backward)
  \item The tibia undergoes coupled \textbf{internal rotation} of approximately
        $15^\circ$ (the ``screw-home'' mechanism) \cite{walker1972quantitative,komistek2003invivo}
  \item The contact area shifts, changing the effective moment arm
\end{enumerate}

This coupled motion means the knee ISA migrates superiorly and posteriorly during
flexion. The four-bar linkage model \cite{OConnor1989} provides a simplified but
insightful mechanical analog:

\begin{lstlisting}[language=Python, caption={Instantaneous screw axis computation for a biological joint.}]
import numpy as np
from typing import Tuple

def compute_isa(
    omega: np.ndarray,
    v: np.ndarray
) -> Tuple[np.ndarray, np.ndarray, float]:
    """Compute the instantaneous screw axis from a twist.

    Parameters
    ----------
    omega : np.ndarray, shape (3,)
        Angular velocity vector.
    v : np.ndarray, shape (3,)
        Linear velocity of the body-fixed frame origin.

    Returns
    -------
    Tuple[np.ndarray, np.ndarray, float]
        (point_on_axis, axis_direction, pitch)
    """
    omega_norm = np.linalg.norm(omega)
    if omega_norm < 1e-10:
        # Pure translation
        direction = v / np.linalg.norm(v)
        return np.zeros(3), direction, np.inf

    direction = omega / omega_norm
    pitch = np.dot(omega, v) / omega_norm**2
    point = np.cross(omega, v) / omega_norm**2

    return point, direction, pitch


def knee_isa_trajectory(
    flexion_angles: np.ndarray,
    roll_slide_ratio: float = 0.6
) -> np.ndarray:
    """Estimate knee ISA position during flexion.

    Simple model: ISA migrates posteriorly and superiorly
    with flexion, governed by the roll-to-slide ratio.

    Parameters
    ----------
    flexion_angles : np.ndarray
        Array of flexion angles in degrees.
    roll_slide_ratio : float
        Ratio of rolling to total contact displacement.
        Pure rolling = 1.0, pure sliding = 0.0. A value of
        approximately 0.6 is consistent with the in-vitro
        tibiofemoral measurements of Walker et al. (1972)
        [walker1972quantitative].

    Returns
    -------
    np.ndarray, shape (N, 3)
        ISA positions in the tibial reference frame.
    """
    # Approximate femoral condyle radius (m); representative
    # adult value from Walker et al. (1972) [walker1972quantitative]
    # and standard knee-anatomy references [nigg2007biomechanics].
    R_condyle = 0.030
    isa_positions = np.zeros((len(flexion_angles), 3))

    for i, theta in enumerate(np.radians(flexion_angles)):
        # Posterior migration due to rolling
        posterior = R_condyle * theta * roll_slide_ratio
        # Superior migration (ISA moves up with flexion)
        superior = R_condyle * (1.0 - np.cos(theta * 0.5))

        isa_positions[i] = [-posterior, superior, 0.0]

    return isa_positions
\end{lstlisting}

\subsection{The Shoulder: The Most Complex Joint}

The glenohumeral joint has the largest range of motion of any joint in the body,
achieved at the expense of stability. The ``ball-and-socket'' description is a
gross simplification: the humeral head translates $2$--$4$~mm during abduction
\cite{harryman1990translation}, and the instantaneous center of rotation shifts
continuously.

The shoulder complex involves four joints:
\begin{enumerate}
  \item \textbf{Glenohumeral} (GH): primary ball-and-socket, 3 DOF rotation + $\sim$3~mm translation \cite{harryman1990translation}
  \item \textbf{Acromioclavicular} (AC): 3 DOF rotation, small range
  \item \textbf{Sternoclavicular} (SC): 3 DOF rotation, small range
  \item \textbf{Scapulothoracic} (ST): not a true synovial joint; the scapula glides on
        the thorax, constrained as a \emph{surface contact} rather than a pin joint
\end{enumerate}

The scapulothoracic ``joint'' is particularly interesting from a mathematical
perspective: it is a \emph{holonomic constraint} that the scapula must remain on
the thoracic surface, effectively acting as a 2-DOF constraint (the scapula can
translate on the surface but cannot lift off).

\section{Ligament Models}

Ligaments constrain joint motion by resisting excessive displacement. They are
modeled as nonlinear springs with:
\begin{equation}
  F_{\text{lig}}(\varepsilon) = \begin{cases}
    0 & \varepsilon \leq 0 \text{ (slack)}, \\
    k_1 \varepsilon^2 / (4\varepsilon_1) & 0 < \varepsilon \leq 2\varepsilon_1 \text{ (toe region)}, \\
    k_2(\varepsilon - \varepsilon_1) & \varepsilon > 2\varepsilon_1 \text{ (linear region)},
  \end{cases}
  \label{eq:ch4:ligament-force}
\end{equation}
where $\varepsilon = (\ell - \ell_0) / \ell_0$ is the strain, $\ell_0$ is the
reference length, $\varepsilon_1 \approx 0.03$ is the transition strain, and
$k_1, k_2$ are stiffness parameters \cite{Blankevoort1991}.

\section{Spinal Mechanics}

The spine presents a unique modeling challenge: 24 vertebrae connected by intervertebral
discs, each joint having 6 coupled degrees of freedom. The functional spinal unit (FSU)
consisting of two adjacent vertebrae and their connecting disc is the basic mechanical
unit.

The intervertebral disc consists of:
\begin{itemize}
  \item \textbf{Nucleus pulposus}: a gel-like center that distributes compressive loads
  \item \textbf{Annulus fibrosus}: concentric rings of collagen fibers that resist
        torsion and bending
\end{itemize}

A common simplified model treats each FSU as a 6-DOF spring:
\begin{equation}
  \bm{F}_{\text{disc}} = K_{\text{disc}} \cdot \bm{\delta},
  \label{eq:ch4:disc-stiffness}
\end{equation}
where $K_{\text{disc}} \in \Reals^{6 \times 6}$ is the stiffness matrix and
$\bm{\delta} \in \Reals^6$ is the generalized displacement (3 translations + 3 rotations).

\section*{Exercises}

\begin{enumerate}
  \item \textbf{ISA Computation.}
        For a knee flexing at $60^\circ/\text{s}$ with $2$~mm/s of posterior translation,
        compute the ISA direction, pitch, and location.

  \item \textbf{Screw-Home Mechanism.}
        Plot the coupled internal rotation of the tibia as a function of knee flexion
        using the model: $\gamma = 15^\circ \cdot (1 - \cos(\theta / 2))$ for
        $\theta \in [0, 90^\circ]$.

  \item \textbf{Shoulder Kinematics.}
        Using screw axis theory, express the configuration of the hand (end-effector)
        as a product of exponentials through the SC, AC, GH, and elbow joints.

  \item \textbf{(Programming.)} Implement the ligament force model and plot the
        force-strain curve. Identify the toe region and the linear region.

  \item \textbf{(Programming.)} Create a 3D visualization of the knee ISA trajectory
        during flexion from $0^\circ$ to $120^\circ$.
\end{enumerate}
